% --- Reinforcement Learning and Post-Training (PPO, DPO) ---
% ============================================================================
\section{Reinforcement Learning and Post-Training}
% Slides below adapted from earlier KAUST Academy lectures (see per-frame comments).
% ============================================================================

% --- Source: Lecture-1_Introduction_Reinforcement_Learning (sections/intro/policy.tex) ---
\begin{frame}{What Is a Policy?}
  \begin{itemize}\itemsep0.4em
    \item[\textcolor{red}{$\blacktriangleright$}] A \textbf{policy} $\pi$ specifies how an agent (or model) chooses actions.
    \item[\textcolor{red}{$\blacktriangleright$}] \textbf{Deterministic:} $\pi(s) = a$.
    \item[\textcolor{red}{$\blacktriangleright$}] \textbf{Stochastic:} $\pi(a \mid s) = \Pr(a_t = a \mid s_t = s)$.
    \item[\textcolor{red}{$\blacktriangleright$}] For LLMs: $s$ is the prompt/context, $a$ is the next token; the policy is $\pi_\theta(\text{token} \mid \text{context})$.
    \item[\textcolor{red}{$\blacktriangleright$}] \textbf{Goal:} find $\pi^\star$ that maximizes expected cumulative reward (or, in RLHF, human-aligned reward).
  \end{itemize}
\end{frame}

% --- Source: Lecture-3_Policy_Grad+REINFORCE+Actor_Critic (sections/policygrad+reinforce+actor/policy_gradient.tex) ---
\begin{frame}{Policy Gradients: Learn the Policy Directly}
  \begin{itemize}\itemsep0.35em
    \item[\textcolor{red}{$\blacktriangleright$}] Value-based methods (Q-learning) can be hard in huge state spaces.
    \item[\textcolor{red}{$\blacktriangleright$}] Often easier to parameterize a policy class $\Pi = \{\pi_\theta \mid \theta \in \mathbb{R}^m\}$.
    \item[\textcolor{red}{$\blacktriangleright$}] Maximize expected return $\mathcal{J}(\theta) = \mathbb{E}\big[\sum_{t \geq 0} \gamma^t r_t \mid \pi_\theta\big]$.
    \item[\textcolor{red}{$\blacktriangleright$}] \textbf{Policy gradient:} update $\theta$ along $\nabla_\theta \mathcal{J}(\theta)$ (gradient ascent on parameters).
    \item[\textcolor{red}{$\blacktriangleright$}] PPO and GRPO are policy-gradient-style updates on language-model weights.
  \end{itemize}
\end{frame}

% --- Source: Lecture-4_Policy_Search (sections/policy-search/ppo.tex) ---
\begin{frame}{Proximal Policy Optimization (PPO)}
    \begin{itemize}\itemsep0.32em
      \item[\textcolor{red}{$\blacktriangleright$}] \textbf{Problem:} large policy updates can collapse performance.
      \item[\textcolor{red}{$\blacktriangleright$}] \textbf{Idea:} take the biggest safe improvement step on the current batch of data.
      \item[\textcolor{red}{$\blacktriangleright$}] \textbf{Clipped objective} (widely used variant): with $r_t(\theta) = \frac{\pi_\theta(a_t \mid s_t)}{\pi_{\theta_k}(a_t \mid s_t)}$,
        \[
          \mathcal{L}^{\mathrm{CLIP}} = \mathbb{E}_t\Big[\min\big(r_t(\theta)\hat{A}_t,\; \mathrm{clip}(r_t(\theta), 1-\epsilon, 1+\epsilon)\hat{A}_t\big)\Big]
        \]
      \item[\textcolor{red}{$\blacktriangleright$}] $\epsilon$ is small (e.g., $0.2$)---limits how far the new policy can drift per update.
    \end{itemize}
  \end{frame}

% --- Source: Lecture-3_Policy_Grad+REINFORCE+Actor_Critic (sections/policygrad+reinforce+actor/actor_critic.tex) ---
% % \begin{frame}{Actor--Critic Methods (Why PPO Needs a Critic)}
% %   \begin{itemize}\itemsep0.35em
% %     \item[\textcolor{red}{$\blacktriangleright$}] \textbf{Actor} = policy $\pi_\theta$ (chooses actions / generates tokens).
% %     \item[\textcolor{red}{$\blacktriangleright$}] \textbf{Critic} = value or advantage estimator (judges how good actions were).
% %     \item[\textcolor{red}{$\blacktriangleright$}] Advantage $\hat{A}_t$ measures whether an action beat expectation---used in PPO updates.
% %     \item[\textcolor{red}{$\blacktriangleright$}] Classic RLHF-PPO keeps \textbf{policy + reference + reward + critic} in memory.
% %     \item[\textcolor{red}{$\blacktriangleright$}] \textbf{GRPO} later drops the critic by using group-relative rewards within a batch of completions.
% %   \end{itemize}
% % \end{frame}

% % --- Source: Lecture-8_Fine_Tuning_LLM+RLHF (sections/llm-rlhf/reinforcement-learning-with-human-feedback-rlhf.tex) ---
% \begin{frame}{RLHF: Why Not Stop at SFT?}
%   \begin{itemize}\itemsep0.4em
%     \item[\textcolor{red}{$\blacktriangleright$}] Supervised fine-tuning (SFT) teaches format and task following, not \textbf{what humans prefer}.
%     \item[\textcolor{red}{$\blacktriangleright$}] \textbf{RLHF} (Reinforcement Learning from Human Feedback) adds reward learning and policy optimization.
%     \item[\textcolor{red}{$\blacktriangleright$}] Standard pipeline (InstructGPT / ChatGPT era): SFT $\rightarrow$ reward model $\rightarrow$ PPO.
%     \item[\textcolor{red}{$\blacktriangleright$}] Reasoning models (R1, o-series) often replace human preference rewards with \textbf{verifiable} rewards (RLVR) and use \textbf{GRPO}.
%   \end{itemize}
% \end{frame}

% \begin{frame}{RL from Human Feedback (RLHF)}
%   \begin{enumerate}\itemsep0.4em
%     \item[\textcolor{red}{$\blacktriangleright$}] \textbf{Supervised fine-tuning (SFT):} Fine-tune a language model using a dataset $D_{\mathrm{SFT}} = \{(x_i^{(0)}, y_i^*)\}_{i=1}^N$.
    
%     \quad \textit{Example:} $x$ is a prompt, $y$ is a human-written response.
%     \item[\textcolor{red}{$\blacktriangleright$}] \textbf{Reward modeling:} Train a reward model $r_\phi(x, y)$ using preference data $D_{\mathrm{RM}} = \{(x_j, y_j^+, y_j^-)\}_{j=1}^M$.
    
%     \quad \textit{Example:} $x$ is a prompt, $y$ is a model-generated response.
%     \item[\textcolor{red}{$\blacktriangleright$}] \textbf{Reinforcement learning (RL):} Further fine-tune the language model from step 1 (call it $p_\theta$) using a reinforcement learning algorithm:
    
%     \[
%     p_\theta = \text{RL}(D_{\mathrm{RM}}, p_\theta^{\mathrm{SFT}})
%     \]
%     \quad \textit{Example:} as in this pipeline, PPO is the RL algorithm.
%   \end{enumerate}
% \end{frame}

% \begin{frame}{RLHF Pipeline}
%   \centering
%   \includegraphics[width=0.92\textwidth,height=0.82\textheight,keepaspectratio]{assets/Lecture-8_Fine_Tuning_RLHF/slide_022.png}\\[-0.2em]
%   {\scriptsize Lecture-8\_Fine\_Tuning\_LLM+RLHF.}
% \end{frame}

% % --- Source: Lecture-8_Fine_Tuning_LLM+RLHF (sections/llm-rlhf/reward-model-rm.tex) ---
% \begin{frame}{Reward Model (RM)}
%   \begin{itemize}\itemsep0.35em
%     \item[\textcolor{red}{$\blacktriangleright$}] Given two responses $y_A$, $y_B$ to the same prompt, humans pick the better one.
%     \item[\textcolor{red}{$\blacktriangleright$}] Train $r_\phi$ so preferred responses score higher:
%       \[
%         L_{\mathrm{RM}} = -\log \sigma\big(r_\phi(y_A) - r_\phi(y_B)\big)
%       \]
%     \item[\textcolor{red}{$\blacktriangleright$}] The RM is a \textbf{proxy} for human judgment during PPO.
%     \item[\textcolor{red}{$\blacktriangleright$}] For reasoning RLVR: the ``reward model'' is often a \textbf{rule} (math checker, unit tests)---no trainable RM needed.
%   \end{itemize}
% \end{frame}

% % --- Source: Lecture-8_Fine_Tuning_LLM+RLHF (sections/llm-rlhf/rlhf-training-proximal-policy-optimization-ppo.tex) ---
% \begin{frame}{PPO for RLHF}
%   \small
%   \begin{itemize}\itemsep0.3em
%     \item[\textcolor{red}{$\blacktriangleright$}] \textbf{Objective:} maximize reward from the RM.
%     \item[\textcolor{red}{$\blacktriangleright$}] \textbf{PPO loss} (same clipped ratio as classic PPO):
%       \[
%         L_{\mathrm{PPO}} = \mathbb{E}\Big[\min\big(r_t(\theta)\hat{A}_t,\; \mathrm{clip}(r_t(\theta), 1-\epsilon, 1+\epsilon)\hat{A}_t\big)\Big]
%       \]
%     \item[\textcolor{red}{$\blacktriangleright$}] \textbf{KL penalty:} $L_{\mathrm{KL}} = \beta\, \mathrm{KL}(\pi_\theta \,\|\, \pi_{\mathrm{SFT}})$ keeps the policy near the SFT checkpoint.
%     \item[\textcolor{red}{$\blacktriangleright$}] \textbf{Pros:} stable; \textbf{cons:} expensive rollouts, sensitive to RM errors, four models in memory.
%   \end{itemize}
% \end{frame}

\begin{frame}{PPO --- Pros and Cons}
  \begin{itemize}\itemsep0.4em
    \item[\textcolor{red}{$\blacktriangleright$}] \textbf{Stable} optimization with clipped updates.
    \item[\textcolor{red}{$\blacktriangleright$}] Encourages exploration of phrasings at each rollout.
    \item[\textcolor{red}{$\blacktriangleright$}] \textbf{Expensive:} many rollouts, critic + RM + reference in memory.
    \item[\textcolor{red}{$\blacktriangleright$}] \textbf{Sensitive} to reward-model mistakes and distribution shift.
  \end{itemize}
\end{frame}

% --- Source: Lecture-8_Fine_Tuning_LLM+RLHF (sections/llm-rlhf/direct-preference-optimization-dpo.tex) ---
\begin{frame}{Direct Preference Optimization (DPO)}
  \begin{itemize}\itemsep0.35em
    \item[\textcolor{red}{$\blacktriangleright$}] \textbf{Bypasses} explicit RM training and RL rollouts (Rafailov et al., 2023).
    \item[\textcolor{red}{$\blacktriangleright$}] Optimize the policy \textbf{directly} from preference pairs $(y^+, y^-)$.
    \item[\textcolor{red}{$\blacktriangleright$}] \textbf{DPO loss} (simplified form):
      \[
        L_{\mathrm{DPO}} = -\log \sigma\Big(\beta \log \tfrac{\pi_\theta(y^+)}{\pi_0(y^+)} - \beta \log \tfrac{\pi_\theta(y^-)}{\pi_0(y^-)}\Big)
      \]
    \item[\textcolor{red}{$\blacktriangleright$}] $\pi_0$ = frozen reference (SFT) policy; $\beta$ = temperature.
    \item[\textcolor{red}{$\blacktriangleright$}] Often \textbf{simpler and more stable} than PPO for alignment---but still preference-based, not verifiable math/code rewards.
  \end{itemize}
\end{frame}

\begin{frame}{Direct Preference Optimization (DPO)}
  \begin{figure}
    \centering
    \includegraphics[width=0.92\textwidth,height=0.82\textheight,keepaspectratio]{assets/Lecture-8_Fine_Tuning_RLHF/slide_031.png}
    \caption*{DPO pipeline: pre-training, supervised fine-tuning, then direct preference optimization from human preference pairs. Source: Rafailov et al. (2023), \href{https://arxiv.org/abs/2305.18290}{arXiv:2305.18290}}
  \end{figure}
\end{frame}

% \begin{frame}{From Alignment RL to Reasoning RL}
%   \small
%   \begin{tabular}{@{}lll@{}}
%     \toprule
%     \textbf{Method} & \textbf{Reward} & \textbf{Typical use} \\
%     \midrule
%     PPO + RM & Human preferences & Chat alignment \\
%     DPO & Preference pairs & Chat alignment (no RL loop) \\
%     GRPO & Verifiable outcome (RLVR) & Math, code, logic reasoning \\
%     GSPO & Same as GRPO; sequence-level clip & Long chains, MoE scale \\
%     \bottomrule
%   \end{tabular}
%   \vspace{0.5em}
%   \begin{itemize}\itemsep0.3em
%     \item[\textcolor{red}{$\blacktriangleright$}] DeepSeek-R1 uses \textbf{GRPO} with rule-based rewards, not classical RLHF.
%     \item[\textcolor{red}{$\blacktriangleright$}] Next: \textbf{GRPO}, then \textbf{GSPO} for stable long-chain and MoE RL.
%   \end{itemize}
% \end{frame}

% \begin{frame}[plain]
%     \centering\vfill
%     {\LARGE\bfseries Part 3: GRPO}\\[0.4em]
%     {\large Group Relative Policy Optimization for reasoning}
%     \vfill
% \end{frame}
